% ============================================================
%  B 组样品制备 SOP（用 09-04 剩余料）— 可交付操作版
%  样品定义：中层葡萄糖 2 g + 空气预氧化 220 °C + Zr 0.64 g + Ar 900 °C
%  编译：在本目录运行  xelatex "B组制备SOP_用剩料.tex"
%  版本：2026-09-22
% ============================================================
\documentclass[a4paper, 11pt]{ctexart}
\usepackage{geometry}
\geometry{left=1.8cm, right=1.8cm, top=2.0cm, bottom=2.0cm}
\usepackage{booktabs}
\usepackage{array}
\usepackage{tabularx}
\usepackage{enumitem}
\usepackage{xcolor}
\setlist[itemize]{itemsep=2pt}
\setlist[enumerate]{itemsep=2pt}
\renewcommand{\arraystretch}{1.25}
\begin{document}

\begin{center}
{\LARGE\bfseries B 组样品制备 SOP（用剩余料，预氧化版）}\\[4pt]
{\normalsize 静电纺丝 SiZrOC 纤维 ｜ 含糖 2\,g ｜ \textbf{空气预氧化 220\,°C} ｜ \ce{Ar} 900\,°C ｜ 2026-09-22 版}
\end{center}

\section*{0. 起点与样品定义}
\begin{itemize}
  \item \textbf{起点（本版关键）}：使用 \textbf{09-04 配制的剩余料}——\textbf{基液}与\textbf{已溶 PVP 的外层液 / 中层液}（均\textbf{尚未滴加正丙醇锆}、未加糖水）。目视状态正常（透明、不拉丝）。
  \item \textbf{放置时间提醒}：该批料配制于 09-04，至本次使用已放置约 18 天。硅烷会持续缩合、DMF 吸湿，因此\textbf{必须先做第 1 步检查与试纺}；若已增稠/凝胶，见 §5 的"重配方案"。
  \item \textbf{样品定义（B 组）}：硅烷（MTMS:DMDMS = 5:1）陶瓷骨架 + PVP 纺丝助剂；中层加\textbf{葡萄糖 2.00\,g}与\textbf{正丙醇锆 0.64\,g}（$\approx$ Zr:Si 5\,mol\%）；三层纺丝（外--中--外）；\textbf{先空气预氧化 220\,°C 成壳}，再 \ce{Ar} 900\,°C 热解 1\,h。
  \item \textbf{用途}：B 是 A（无预氧化）的对照，用来回答"\textbf{预氧化能否做出致密壳/表面突起}"。
\end{itemize}

\section*{1. 检查与取料（第一步，必做）}
\begin{enumerate}
  \item \textbf{目视检查}每个烧杯：应为透明、无浑浊、无白色絮状、无凝胶块；
  \item \textbf{挑丝检查}：玻璃棒挑起看拉丝长度、挂壁厚薄、气泡悬浮时间，写进记录表；
  \item \textbf{小样试纺}：取 2--3\,mL 液体试纺，确认\textbf{射流稳定}（泰勒锥稳定、无堵针/滴液）才继续；不稳则调电压/喷射距离/挤出速率；
  \item \textbf{取料量}：
    \begin{itemize}
      \item 外层液：约 \textbf{23.4\,g}（供两个外层各 10\,mL + 试纺 1--3\,mL）；
      \item 中层液：约 \textbf{12.4\,g}（供中层 10\,mL + 试纺）；
      \item 若剩料量不足，可按 §5 新配补足。
    \end{itemize}
\end{enumerate}

\section*{2. 加料（顺序不可颠倒）}
\begin{table}[htbp]
\centering
\caption{本次需补加的组分（其余用量沿用剩料）}
\begin{tabularx}{\textwidth}{l l X}
\toprule
步骤 & 组分与用量 & 操作要点 \\
\midrule
1 & 正丙醇锆 \textbf{0.64\,g} & 注射器\textbf{缓慢滴入中层液}，边加边搅；搅 \textbf{30\,min}（不必延长） \\
2 & d-无水葡萄糖 \textbf{2.00\,g} + 去离子水 \textbf{3.0\,mL} & 40--50\,°C 搅 5--10\,min 溶清，冷却至 $<$40\,°C；\textbf{现配现用} \\
3 & 上述糖水全部 & 加入中层液，搅 \textbf{15--30\,min} $\rightarrow$ \textbf{30\,min 内必须开纺} \\
\bottomrule
\end{tabularx}
\end{table}

\noindent\textbf{关键顺序}：先锆、后糖水。锆遇游离水易水解析出，所以糖水（带 3\,mL 水）必须最后加。

\section*{3. 纺丝}
\begin{enumerate}
  \item 三个 10\,mL 注射器分别灌装：\textbf{外层液 $\times$2、中层液 $\times$1}，装\textbf{全金属针头}（塑料件会被 DMF 溶胀）；
  \item \textbf{顺序（外--中--外）}：外层 10\,mL $\rightarrow$ 中层 10\,mL $\rightarrow$ 外层 10\,mL；
  \item 挤出速率 \textbf{1.5\,mL/h}；每层开始前试纺至射流稳定（约 5--15\,min，不按体积计）；
  \item \textbf{三层连续不中断}：每层约 6.7\,h，合计约 \textbf{20\,h（跨夜）}；
  \item 纺完连同硅油纸取下，称\textbf{初重}，室温平放晾置（防折叠、防沾尘、防潮）；
  \item 记录：纺丝起止时间、环境状况（有湿度计则记数值）、液体定性黏度变化。
\end{enumerate}

\section*{4. 热处理（两步，必须分开）}
\subsection*{4.1 第一步：空气预氧化成壳（马弗炉）}
\begin{table}[htbp]
\centering
\caption{空气预氧化程序（B 组关键步骤）}
\begin{tabularx}{\textwidth}{l l X}
\toprule
段 & 参数 & 说明 \\
\midrule
升温 1 & 室温 $\rightarrow$ \textbf{1--2\,°C/min} $\rightarrow$ 100\,°C，保温 \textbf{30\,min} & 先赶掉残余 DMF/水，避免冲击 \\
升温 2 & 100 $\rightarrow$ 150\,°C，保温 \textbf{30\,min} & \textbf{新增过渡段}：葡萄糖熔点 146\,°C，先熔融/脱水，避免毛细流动板结 \\
升温 3 & 150 $\rightarrow$ \textbf{220\,°C}，保温 \textbf{60\,min} & 全程\textbf{不超过 220\,°C}（PVP 耐氧上限较低） \\
降温 & 随炉自然降温至 \textbf{$<$60\,°C} 取出 & 低于 60\,°C 再取，避免热冲击 \\
\bottomrule
\end{tabularx}
\end{table}
\noindent 记录\textbf{预氧化前后质量}（可算出预氧化失重率）与\textbf{颜色变化}（应为均匀金黄/浅褐色）。操作时平铺入炉、不折叠。

\subsection*{4.2 第二步：\ce{Ar} 气氛热解（管式炉）}
\begin{enumerate}
  \item \textbf{必须先降温再换气氛}：预氧化样冷却至 $<$60\,°C 后再转移入管式炉，平铺入石英舟（记录舟位编号）；
  \item \ce{Ar} 吹扫 \textbf{30\,min} 后升温；
  \item 程序：室温 $\rightarrow$（\textbf{200\,min}）$\rightarrow$ 200\,°C，保温 \textbf{60\,min}（继续除残余溶剂）$\rightarrow$（\textbf{400\,min}）$\rightarrow$ 900\,°C，保温 \textbf{60\,min}；
  \item \textbf{自然炉冷}，$<$100\,°C 再停气取出；
  \item 称\textbf{终重}、算失重率，装袋贴标签（B + 日期 + 操作者）。
\end{enumerate}

\section*{5. 异常处置与重配方案}
\begin{itemize}
  \item \textbf{剩料已增稠/凝胶（拉长丝不断、挂壁厚、有胶块或浑浊）} $\rightarrow$ \textbf{不要硬纺}，按下面新配：
    \begin{table}[htbp]
    \centering
    \caption{单样品新配方案（剩料不可用时）}
    \begin{tabularx}{\textwidth}{l l X}
    \toprule
    液体 & 用量 & 备注 \\
    \midrule
    硅烷原液 & MTMS \textbf{16.67\,g} + DMDMS \textbf{3.33\,g} & 加去离子水 \textbf{0.2--0.4\,mL}（不加酸），密封磁搅\textbf{过夜} \\
    基液 & 硅烷原液 20.00\,g + DMF \textbf{20.00\,g} & 搅 30\,min，得 40\,g（1:1） \\
    外层液 & 基液 \textbf{20.00\,g} + PVP \textbf{3.44\,g} & 搅 6\,h \\
    中层液 & 基液 \textbf{10.00\,g} + PVP \textbf{1.72\,g} & 搅 6\,h，第 4--5\,h 滴入 0.64\,g 正丙醇锆 \\
    糖水 & 葡萄糖 \textbf{2.00\,g} + 去离子水 \textbf{3.0\,mL} & 现配现用 \\
    \bottomrule
    \end{tabularx}
    \end{table}
  \item \textbf{加锆后出现浑浊/白色絮状} $\rightarrow$ 锆在析出，\textbf{立即停止}并记录，不要纺；
  \item \textbf{纺丝中断丝/滴液频繁} $\rightarrow$ 检查堵针与挤出速率，必要时重灌；
  \item \textbf{预氧化后颜色发黑/发脆} $\rightarrow$ 说明温度过高或时间过长（PVP 被氧化烧损），记录并下次降为 200\,°C 保温 60\,min；
  \item \textbf{高温后膜碎裂} $\rightarrow$ 预氧化不足或升温过快，记录实际曲线。
\end{itemize}

\section*{6. 记录表（操作者填写）}
\begin{table}[htbp]
\centering
\caption{本次操作记录}
\begin{tabularx}{\textwidth}{l l l X}
\toprule
项目 & 计划值 & 实际值 & 现象/备注 \\
\midrule
剩料放置天数 & 约 18 天（09-04 配） & & 目视/挑丝状态 \\
外层液取用量 & $\approx$23.4\,g & & \\
中层液取用量 & $\approx$12.4\,g & & \\
正丙醇锆 & 0.64\,g & & 滴加后是否浑浊 \\
葡萄糖 / 水 & 2.00\,g / 3.0\,mL & & 糖水是否澄清 \\
开纺前糖水加入时刻 & & & 是否 30\,min 内开纺 \\
纺丝起止时间 & & & 是否连续（$\sim$20\,h）\\
初重 / 预氧化后重 & & & 预氧化失重率 \\
预氧化颜色 & 金黄/浅褐 & & \\
终重（热解后） & & & 总失重率 \\
热解程序 & 200\,min $\rightarrow$ 200\,°C 保温 1\,h；400\,min $\rightarrow$ 900\,°C 保温 1\,h & & 舟位编号 \\
\bottomrule
\end{tabularx}
\end{table}

\section*{7. 交付物与后续对比要求}
\begin{enumerate}
  \item 烧结后样品（装袋、编号 B + 日期）；完整记录表；外观照片；
  \item \textbf{SEM 送样}：表面形貌（10k--20k）+ \textbf{截面}（验证壳-芯）；
  \item \textbf{对比规则（与 A 样对比时）}：孔隙率/孔径\textbf{必须同倍率}比较；直径用距离变换中位数 + IQR（不用均值）；条纹间距取 20k/15k/10k 三档平均；
  \item 若条件允许，加做 \textbf{EDS 线扫}看 Zr 是否向表面富集（预氧化成壳的判据之一）。
\end{enumerate}

\end{document}
